% Versión breve basada en el documento compartido por Fernanda.
% Autocontenida: no necesita imágenes ni otros archivos para compilar.
\documentclass[12pt,letterpaper]{article}
\usepackage[letterpaper,margin=1in]{geometry}
\usepackage{iftex}
\ifPDFTeX
  \usepackage[utf8]{inputenc}
\fi
\usepackage[T1]{fontenc}
\usepackage[spanish,es-tabla]{babel}
\usepackage{mathptmx}
\usepackage{graphicx}
\usepackage[section]{placeins}
\usepackage{amsmath}
\usepackage{microtype}
\usepackage{tikz}
\usetikzlibrary{arrows.meta}
\definecolor{scAzul}{HTML}{244C65}
\definecolor{scVerde}{HTML}{32685D}
\definecolor{scGris}{HTML}{F3F5F7}
\usepackage{xurl}
\usepackage[hidelinks]{hyperref}
\setlength{\emergencystretch}{2em}
\hypersetup{pdftitle={SmartCity Analytics: documentación técnica del Sprint 1},
  pdfauthor={Daniel Burger, José Pablo Alpízar, Fernanda Vargas, Kembly Garro}}

\begin{document}
\hypersetup{pageanchor=false}
\begin{titlepage}
    \centering
    \vspace*{1.8cm}
    {\large Instituto Tecnológico de Costa Rica\par}
    \vspace{0.8cm}
    {\large Escuela de Ingeniería en Computadores\par}
    \vspace{0.8cm}
    {\large CE-3191 Base de Datos\par}
    \vspace{1.8cm}
    {\large SmartCity Analytics\par}
    \vspace{0.5cm}
    {\Large\scshape Documentación técnica --- Sprint 1\par}
    \vspace{1.8cm}
    {\large Integrantes:\par}
    \vspace{0.5cm}
    {\large Daniel Burger --- 2024145222\par}
    \vspace{0.2cm}
    {\large José Pablo Alpízar --- 2024088025\par}
    \vspace{0.2cm}
    {\large Fernanda Vargas --- 2022008164\par}
    \vspace{0.2cm}
    {\large Kembly Garro --- 2024134020\par}
    \vfill
    {\large Cartago, Costa Rica\par}
    \vspace{0.2cm}
    {\large Octubre, 2026\par}
\end{titlepage}
\pagenumbering{arabic}
\hypersetup{pageanchor=true}

\section{Descripción del problema}

Las ciudades generan constantemente grandes volúmenes de información provenientes
de distintas fuentes, como sensores, dispositivos IoT y plataformas de monitoreo
urbano. Estos datos son útiles para la gestión municipal, por ejemplo, en calidad
del aire, condiciones climáticas y tráfico vehicular. El reto consiste en gestionar
información con estructuras y necesidades diferentes mediante tecnologías de
almacenamiento complementarias.

Por un lado, existen datos estructurados y relacionales, como ubicaciones,
usuarios y organismos responsables, que requieren consistencia e integridad
referencial. Por otro, las lecturas y los eventos pueden contener atributos
variables y estructuras anidadas. Para abordar esta problemática, el primer sprint
implementa una arquitectura híbrida: PostgreSQL organiza la información mediante
un modelo normalizado hasta la tercera forma normal, mientras que MongoDB permite
trabajar con documentos BSON y Aggregation Framework.

\section{Fuentes de datos}

Tras evaluar distintas opciones, se seleccionaron los datos abiertos del
Ayuntamiento de Madrid. La generación v2 validada contiene \textbf{748.889 lecturas
de enero de 2025}: 91.959 de calidad del aire y 656.930 de tráfico vehicular.
Meteorología fue investigada, pero no forma parte del lote aceptado. El tráfico
corresponde a una selección de puntos de medida, no a todo el conjunto de Madrid.

Como complemento, se utilizan 14.388 accidentes de 2026, agrupados por expediente
y anonimizados en la salida preparada. Estos incidentes se mantienen separados
de las lecturas de 2025 y no se suman a su volumen.

Las fuentes oficiales omiten algunos atributos administrativos solicitados por el
proyecto. Por ello se incorporan 5.133 sensores, 5.133 usuarios y un organismo
\textbf{sintéticos}, generados de forma determinista y con su procedencia
identificada. Estos datos permiten demostrar las entidades del modelo, pero no
acreditan dispositivos físicos, instalaciones ni responsables municipales reales.
Las mediciones observadas no se modifican para completar esos atributos.

\section{Arquitectura}

La arquitectura de SmartCity Analytics se concibe bajo un enfoque híbrido y
modular para atender las diferencias entre datos relacionales y documentales.

\subsection{Enfoque híbrido de persistencia}
\begin{itemize}
    \item \textbf{Capa relacional:} PostgreSQL almacena catálogos, entidades
    administrativas y observaciones mediante catorce tablas normalizadas.
    Las claves, restricciones y triggers comprueban las relaciones del modelo.
    \item \textbf{Capa documental:} MongoDB almacena lecturas, incidentes,
    eventos derivados, alertas y entidades administrativas. Los documentos
    conservan metadatos de calidad, fechas y procedencia.
\end{itemize}

\subsection{Flujo de integración e ingestión}
\begin{enumerate}
    \item \textbf{Extracción y normalización:} Se procesan archivos históricos
    para revisar fechas, unidades y referencias a los originales mediante SHA-256.
    \item \textbf{Carga y control de integridad:} Se cargan las lecturas aceptadas
    y se comprueba que repetir el procesamiento no duplique ni altere su contenido.
    \item \textbf{Consulta y verificación:} Se ejecutan las consultas de ambos
    motores y se contrastan los resultados con cálculos independientes en Python.
\end{enumerate}

Ambos motores reciben la misma generación preparada; no se implementa una
sincronización continua ni una transacción distribuida entre ellos.

\begin{figure}[htbp]
\centering
\resizebox{\linewidth}{!}{%
\begin{tikzpicture}[
 bloque/.style={draw=scAzul!65,rounded corners=3pt,fill=scGris,
 align=flush center,text width=3.7cm,minimum height=1.15cm,
 inner sep=6pt,font=\small},
 flecha/.style={-{Latex[length=2mm]},draw=scAzul,line width=0.8pt}]
\node[bloque] (origen) at (0,0)
 {\textbf{1. Fuentes oficiales}\\Archivos históricos\\Aire y tráfico, enero 2025};
\node[bloque,fill=scAzul!6] (preparacion) at (4.6,0)
 {\textbf{2. Preparación}\\Fechas, unidades y calidad\\Referencias y hashes};
\node[bloque,draw=scVerde,fill=scVerde!8] (lote) at (9.2,0)
 {\textbf{3. Paquete aceptado}\\748.889 lecturas\\misma generación v2};
\node[bloque,fill=scAzul!6] (pg) at (6.7,-2.3)
 {\textbf{PostgreSQL}\\Catorce tablas en 3FN\\Relaciones y vistas};
\node[bloque,draw=scVerde,fill=scVerde!8] (mongo) at (11.7,-2.3)
 {\textbf{MongoDB}\\Documentos BSON\\Índices y agregaciones};
\node[bloque,text width=6.8cm] (comparacion) at (9.2,-4.5)
 {\textbf{4. Verificación independiente}\\Consultas frente a Python; identidad, contenido y replay};
\draw[flecha] (origen) -- (preparacion);
\draw[flecha] (preparacion) -- (lote);
\coordinate (reparto) at (9.2,-1.15);
\draw[draw=scAzul,line width=0.8pt] (lote.south) -- (reparto);
\fill[scAzul] (reparto) circle (1.3pt);
\draw[flecha] (reparto) -| (pg.north);
\draw[flecha] (reparto) -| (mongo.north);
\draw[flecha] (pg.south) -- ++(0,-0.35) -| ([xshift=-1.6cm]comparacion.north);
\draw[flecha] (mongo.south) -- ++(0,-0.35) -| ([xshift=1.6cm]comparacion.north);
\end{tikzpicture}%
}
\caption{Flujo del lote observado. Los dos motores reciben el mismo paquete;
las flechas no representan sincronización entre bases ni procesamiento distribuido.}
\label{fig:flujo}
\end{figure}

\section{Justificación tecnológica}

\subsection{PostgreSQL y tercera forma normal}
PostgreSQL permite representar relaciones mediante claves primarias y foráneas,
evitando referencias inconsistentes. La normalización separa información de
fuentes, estaciones, tipos de dispositivos y observaciones para reducir redundancias
y anomalías de actualización. Las consultas SQL permiten relacionar estos datos
y obtener resultados verificables.

\subsection{MongoDB}
MongoDB permite conservar documentos con estructuras anidadas y atributos
variables. Sus validadores establecen controles sobre las colecciones principales,
mientras que Aggregation Framework facilita filtros, agrupaciones, conteos y
promedios. Esta flexibilidad resulta útil para representar lecturas y documentos
con distinta estructura, sin implicar ausencia de reglas de calidad.

\subsection{Sinergia del enfoque híbrido}
Ambos paradigmas permiten estudiar representaciones complementarias de los datos.
La consistencia entre las lecturas se comprueba comparando identificadores y
hashes de contenido. No se realizó una comparación de rendimiento que permita
declarar un motor superior al otro.

\section{Modelo relacional}

La implementación PostgreSQL se diseñó bajo un modelo de catorce tablas normalizado
hasta 3FN respecto de las dependencias funcionales documentadas. Los campos JSONB
se consideran dominios SQL indivisibles; no se afirma normalización recursiva de
sus propiedades. El modelo organiza la información de forma que las observaciones
no repitan atributos determinados por los catálogos.

Se distingue entre estación, canal lógico y sensor físico. Una estación identifica
un sitio de la fuente; un canal representa una serie de una magnitud; una
observación registra un valor en un instante. La relación entre canal y sensor
solo se incorpora cuando se conoce. En las lecturas públicas actuales no existen
asignaciones físicas acreditadas.

\begin{figure}[htbp]
\centering
{\small\textbf{A. Observaciones y catálogos}}\par\smallskip
\resizebox{\linewidth}{!}{%
\begin{tikzpicture}[
 entidad/.style={draw=scAzul!65,rounded corners=3pt,fill=scAzul!5,
 align=flush center,text width=3.9cm,minimum height=1.2cm,inner sep=5pt,font=\small},
 relacion/.style={-{Latex[length=2mm]},draw=scAzul,line width=0.7pt}]
\node[entidad] (sta) at (0,0)
 {\textbf{stations}\\PK: station\_id\\FK: location\_id $\to$ B};
\node[entidad] (cha) at (5.1,0)
 {\textbf{channels}\\PK: channel\_id\\Serie de una magnitud};
\node[entidad] (obs) at (10.2,0)
 {\textbf{observations}\\PK: measurement\_id\\Valor, instante y procedencia};
\node[entidad] (src) at (0,-2.1)
 {\textbf{sources}\\PK: source\_id};
\node[entidad] (mag) at (5.1,-2.1)
 {\textbf{measurement\_types}\\PK: measurement\_type\\Magnitud y unidad};
\node[entidad] (ver) at (10.2,-2.1)
 {\textbf{source\_versions}\\PK: source\_version\\Snapshot original};
\draw[relacion] (cha.west) -- (sta.east);
\draw[relacion] (obs.west) -- (cha.east);
\draw[relacion] (sta.south) -- (src.north);
\draw[relacion] (cha.south) -- (mag.north);
\draw[relacion] (obs.south) -- (ver.north);
\draw[relacion] (mag.west) -- (src.east);
\draw[relacion] (ver.south) -- ++(0,-0.55) -| (src.south);
\end{tikzpicture}%
}
\par\medskip
{\small\textbf{B. Territorio y administración}}\par\smallskip
\resizebox{\linewidth}{!}{%
\begin{tikzpicture}[
 entidad/.style={draw=scVerde!70,rounded corners=3pt,fill=scVerde!5,
 align=flush center,text width=3.9cm,minimum height=1.2cm,inner sep=5pt,font=\small},
 relacion/.style={-{Latex[length=2mm]},draw=scVerde,line width=0.7pt},
 nota/.style={font=\scriptsize,align=flush center}]
\node[entidad] (cat) at (0,0)
 {\textbf{sensor\_categories}\\PK: category\_id};
\node[entidad] (typ) at (5.1,0)
 {\textbf{sensor\_types}\\PK: sensor\_type};
\node[entidad] (sen) at (10.2,0)
 {\textbf{sensors}\\PK: sensor\_id\\FK: station\_id $\to$ A};
\node[entidad] (loc) at (0,-2.4)
 {\textbf{locations}\\PK: location\_id\\FK: inventory\_version $\to$ A};
\node[entidad] (lnk) at (5.1,-2.4)
 {\textbf{channel\_sensors}\\PK/FK: channel\_id $\to$ A\\Asignación física conocida};
\node[entidad] (org) at (10.2,-2.4)
 {\textbf{organizations}\\PK: organization\_id};
\node[entidad] (zon) at (0,-5.1)
 {\textbf{zones}\\PK: zone\_id\\FK propia: parent\_zone\_id};
\node[entidad] (usr) at (10.2,-5.1)
 {\textbf{users}\\PK: user\_id\\Sin FK en este esquema};
\draw[relacion] (typ.west) -- (cat.east);
\draw[relacion] (sen.west) -- (typ.east);
\draw[relacion] (sen.south) -- (org.north);
\draw[relacion] (lnk.north) -- ++(0,0.35) -| ([xshift=-0.75cm]sen.south);
\draw[relacion] (loc.south) -- (zon.north);
\draw[relacion] (zon.west) -- ++(-0.45,0) |- ([yshift=-0.25cm]zon.south) -- (zon.south);
\node[nota,text width=4cm] at (5.1,-5.1)
 {A y B son dos vistas del mismo esquema.\\Las referencias entre paneles se indican en los campos FK.};
\end{tikzpicture}%
}
\caption{Las catorce tablas, agrupadas para facilitar la lectura. Las flechas
apuntan de hija a padre; las cuatro FK entre paneles se indican dentro de las cajas.
PK: clave primaria. FK: clave foránea.}
\label{fig:relacional}
\end{figure}

El esquema contiene diecisiete claves foráneas. Cada referencia obligatoria exige
un padre y las opcionales permiten su ausencia. Un padre puede tener varias filas
hijas, salvo la asignación de canal a sensor, limitada a una fila por canal.
Las vistas reconstruyen los atributos compartidos sin duplicarlos en las tablas.
El DDL completo, con atributos y restricciones, está en
\path{sprint1/postgres/schema_v2_3nf.sql}.

\section{Modelo documental}

Para gestionar la información semiestructurada se implementa un entorno documental
basado en MongoDB. La generación v2 utiliza siete colecciones:
\texttt{readings}, \texttt{derived\_events}, \texttt{alerts},
\texttt{incidents}, \texttt{organizations}, \texttt{users} y \texttt{sensors}.
Las lecturas incluyen valor, unidad, calidad, fecha BSON y referencias al origen.

\begin{figure}[htbp]
\centering
\resizebox{\linewidth}{!}{%
\begin{tikzpicture}[
 coleccion/.style={draw=scAzul!65,rounded corners=3pt,align=flush center,
 text width=4.1cm,minimum height=1.2cm,inner sep=5pt,font=\small},
 derivacion/.style={-{Latex[length=2mm]},draw=scAzul,line width=0.8pt},
 referencia/.style={-{Latex[length=2mm]},draw=scVerde,dashed,line width=0.8pt},
 grupo/.style={anchor=south west,font=\small\bfseries,text=scAzul}]
\node[grupo] at (-2.2,1.2) {Lecturas observadas y derivados --- enero de 2025};
\node[coleccion,fill=scAzul!6] (lecturas) at (0,0)
 {\textbf{readings}\\748.889 lecturas\\Valor, fecha, canal y origen};
\node[coleccion,fill=scAzul!6] (eventos) at (5.2,0)
 {\textbf{derived\_events}\\10.841 eventos demo\\Intensidad $>2\,000$ veh/h};
\node[coleccion,fill=scAzul!6] (alertas) at (10.4,0)
 {\textbf{alerts}\\10.841 alertas locales\\Replay; no avisos oficiales};
\draw[derivacion] (lecturas) -- (eventos);
\draw[derivacion] (eventos) -- (alertas);
\node[grupo,text=scVerde] at (-2.2,-1.9) {Administración sintética --- sin enlace a lecturas públicas};
\node[coleccion,draw=scVerde,fill=scVerde!5] (organismos) at (0,-2.8)
 {\textbf{organizations}\\1 organismo demo};
\node[coleccion,draw=scVerde,fill=scVerde!5] (sensores) at (5.2,-2.8)
 {\textbf{sensors}\\5.133 sensores demo};
\node[coleccion,draw=scVerde,fill=scVerde!5] (usuarios) at (10.4,-2.8)
 {\textbf{users}\\5.133 usuarios demo};
\draw[referencia] (sensores) -- (organismos);
\node[grupo] at (-2.2,-4.15) {Incidentes --- período independiente};
\node[coleccion,fill=scAzul!6] (incidentes) at (0,-5.15)
 {\textbf{incidents}\\14.388 accidentes\\Anonimizados; año 2026};
\node[align=flush left,text width=9.5cm,font=\small,anchor=west] at (2.9,-5.15)
 {\textbf{Flecha continua:} derivación de eventos y alertas.\\
 \textbf{Flecha discontinua:} referencia lógica al organismo.\\
 Las referencias no son claves foráneas MongoDB.};
\end{tikzpicture}%
}
\caption{Las siete colecciones, separadas por origen y período. Las cuatro cajas
azules tienen validador; las tres administrativas no. El laboratorio sintético
se mantiene en otra base.}
\label{fig:documental}
\end{figure}

\subsection{Validación e integridad}
El cargador \path{sprint1/basico/cargar_mongodb_v2.js} define validadores para
lecturas, incidentes, eventos y alertas. Las tres colecciones administrativas
conservan su procedencia sintética, pero no tienen validadores adicionales.
La identidad de cada lectura y el hash de su contenido permiten comprobar que
el replay mantenga los documentos sin sobrescribir cambios incompatibles.

\subsection{Indexación y consultas}
Se registraron trece índices entre las siete colecciones. Las lecturas incorporan
un índice único por canal, fecha y rol temporal, además de índices compuestos
por fuente, magnitud y fecha, y por zona y magnitud. Los eventos, alertas e
incidentes tienen índices orientados a sus consultas. Se capturaron planes de
ejecución, sin afirmar que todos los filtros estén completamente optimizados.

El núcleo complementario de práctica utiliza \texttt{stageReading},
\texttt{staging\_key} y cuarentena para demostrar duplicados y conflictos con
ejemplos acotados. Es una implementación distinta de la carga masiva v2.
Un laboratorio sintético separado demuestra consultas por sensor, por tipo de
dispositivo y recuperación de texto narrativo con atributos variables y revisiones.

\section{Estrategia Big Data}

La estrategia de manejo de volumen se fundamenta en cargas históricas por lotes,
trazabilidad e índices para apoyar las consultas. Las 748.889 lecturas aceptadas
superan los objetivos de 100.000 registros PostgreSQL y 500.000 documentos MongoDB.
Los catálogos, replay, incidentes y datos sintéticos no se suman a ese volumen.

MongoDB carga lotes de 1.000 documentos para evitar mantener todo el archivo en
memoria. Las agregaciones permiten producir métricas y eventos derivados sin
mezclar unidades ni fuentes. Los hashes ayudan a verificar que ambos motores
conserven las mismas identidades y contenidos.

El alcance implementado no incluye recepción en tiempo real ni almacenamiento
distribuido. MongoDB se probó en modo standalone: no se demostraron replica set
ni sharding. Hadoop/HDFS y Spark corresponden a la continuación del proyecto.
El requisito de almacenamiento distribuido MongoDB sigue abierto y su modalidad
suficiente debe aclararse con el docente. Tampoco se implementaron vistas
materializadas ni particionamiento PostgreSQL; su exigencia para este avance
permanece por confirmar. Un índice parcial no equivale a particionamiento.

\section{Procesamiento}

\subsection{Preparación y capa relacional}
El módulo \texttt{preparar\_v2.py} revisa los archivos originales, estandariza
fechas y unidades y aplica reglas de calidad por fuente. Los rechazos se registran
por su motivo; los valores ausentes no se convierten en cero ni se inventan
atributos físicos. Los accidentes se deduplican por expediente y se anonimizan
en la salida, manteniendo separado su período de 2026.

En PostgreSQL, el ejecutor aplica el modelo normalizado en una instancia temporal
y carga catálogos y observaciones. La vista \texttt{observation\_contract}
reconstruye la representación compartida. Se compara con los documentos de
entrada para comprobar que la normalización conserva identificadores, valores,
fechas y metadatos. Las consultas se contrastan con Python.

\subsection{Pipeline documental}
La carga v2 de MongoDB inserta lecturas por lotes y comprueba conteos y replay.
Después deriva eventos de tráfico con intensidad superior a 2.000 veh/h y
genera alertas locales para zonas conocidas. La regla es de demostración:
no se presenta como umbral legal ni como aviso oficial emitido en 2025.

Las consultas emplean filtros, proyecciones, ordenamiento y Aggregation Framework.
Se calculan cantidades, extremos y promedios por grupos compatibles de fuente,
magnitud y unidad. Los resultados se comparan con cálculos independientes sobre
los archivos preparados. Las pruebas temporales no acreditan el estado de un
servidor compartido ni una operación de producción.

\section{Resultados}

\subsection{Modelo relacional}
La carga normalizada PostgreSQL terminó en \texttt{PASS}, con
\textbf{748.889 lecturas aceptadas}. Se registraron catorce tablas, diecisiete FK
y veinticinco índices. La reconstrucción contractual presentó cero diferencias
respecto a la entrada y el replay conservó el lote.

Las diez consultas \texttt{PG01--PG10} se compararon completamente con Python:
\begin{itemize}
    \item \textbf{Catálogos e inventario:} PG01 produjo doce grupos de tipos
    administrativos de dispositivos, no doce magnitudes de aire. PG02 y PG05
    produjeron 5.133 filas; PG04, un organismo demo; PG08, veintiún grupos de zona.
    \item \textbf{Estados y fechas:} PG03 produjo dos estados administrativos
    y PG06 cinco grupos anuales de instalación, con datos sintéticos identificados.
    \item \textbf{Actividad, vista y función:} PG09 produjo 360 canales y PG10
    5.133 filas de conteo por estación. PG07 devolvió cero: al no existir vínculos
    físicos acreditados sensor--canal, ese resultado no demuestra actividad ni
    inactividad de dispositivos. Su lógica se comprobó con casos controlados.
\end{itemize}

\subsection{Modelo documental}
La ejecución masiva MongoDB también terminó en \texttt{PASS}, con 748.889 lecturas:
91.959 de \texttt{MAD-AIR} y 656.930 de \texttt{MAD-TRAFFIC}. Se verificaron las
diez consultas, el replay, siete operaciones CRUD y la reconciliación de
identificadores y hashes con PostgreSQL.

\begin{itemize}
    \item \textbf{Consultas y derivados:} MG02 encontró 748.765 lecturas dentro
    de su ventana temporal; las otras 124 quedan fuera del filtro, no son rechazos.
    Se generaron 10.841 eventos y 10.841 alertas demo, además de almacenar por
    separado 14.388 incidentes anonimizados de 2026.
    \item \textbf{CRUD y laboratorio:} Las siete operaciones se ejecutaron en
    una colección de práctica aislada sin modificar las lecturas fuente.
    El laboratorio sintético verificó dieciocho documentos y trece rechazos,
    con consultas por sensor, por tipo y por texto narrativo.
    \item \textbf{Rendimiento observado:} MG10 examinó 656.930 documentos y
    registró 6.470 ms en la ejecución conservada. Esto no demuestra optimización
    completa ni permite afirmar superioridad de rendimiento entre motores.
\end{itemize}

Los resultados masivos corresponden a ejecuciones identificadas por hashes y
versiones. La suite offline posterior aprobó 66 pruebas Python y 19 Node.
La integración de práctica aprobó respaldo/restauración en ambos motores y una
carrera PostgreSQL de ocho conexiones. Estas pruebas acotadas no sustituyen
las evidencias masivas ni demuestran concurrencia de producción.

\section{Conclusiones}
\begin{enumerate}
    \item \textbf{Complementariedad del modelo híbrido:} PostgreSQL permitió
    organizar las relaciones y restricciones del modelo normalizado, mientras
    que MongoDB permitió consultar documentos con estructuras variables.
    La comparación de identidades y hashes confirmó la consistencia del lote
    observado entre ambos motores.
    \item \textbf{Trazabilidad y control:} Las referencias a originales,
    validaciones y comprobaciones de replay permiten detectar diferencias y
    preservar el contenido aceptado. Los datos administrativos, el laboratorio
    sintético y los eventos derivados se identifican por separado para no
    atribuirles carácter de mediciones o avisos oficiales.
    \item \textbf{Alcance del avance:} Se cuenta con evidencia técnica de
    volumen, consultas, CRUD e integridad dentro de los entornos probados.
    Persisten límites de identidad física, vigencia geográfica y optimización.
    El almacenamiento distribuido MongoDB permanece sin demostración, por lo
    que los resultados no equivalen a aceptación académica global del avance.
\end{enumerate}

\section*{Fuentes y evidencias de consulta}
\begin{itemize}
    \item Ayuntamiento de Madrid: calidad del aire, tráfico y accidentes.
    \url{https://datos.madrid.es/}
    \item Código, modelos, consultas y documentación:
    \url{https://github.com/Alpizar28/SmartCity}
    \item Evidencia PostgreSQL normalizada:
    \path{sprint1/postgres/evidencias/3nf/20261028T071812-55791345.json}.
    Evidencia MongoDB masiva:
    \path{research/datasets/results/expanded-v2-2025-01/execution-d26e40706dcd.json}.
    Los cierres de cada motor describen sus resultados y limitaciones.
\end{itemize}

\end{document}
